\begin{frame}{Ожидаемый вклад и границы вывода}
\begin{enumerate}
\item Способ сопряжения предобученного временного модуля с отдельной биофизической моделью человеческого ECAP.
\item Генератор нейродинамических данных с воспроизводимыми сценариями и проверяемым происхождением.
\item Сравнительная оценка структуры, динамики и предобучения с независимой проверкой по участникам.
\item Условия применимости, чувствительность и оценки неопределённости для численного исследования измерительного канала.
\end{enumerate}
\SlideGap
\textbf{Граница вывода.} Ноцицепция Drosophila не тождественна субъективной боли человека; ECAP не является прямой мерой боли или эффективности лечения.
\SlideGap
Прогноз клинических исходов и настройка стимуляции потребуют самостоятельных разрешённых данных и проверки.
\SlideGap
Наличие и величина дополнительного вклада переноса будут установлены сравнительной проверкой.
\end{frame}
